% 06_interconnect.tex -- AXI channels + handshake, protocol family, crossbar, NoC router.
% Build: tectonic 06_interconnect.tex
\documentclass[tikz,border=10pt,12pt]{standalone}
% soc_style.tex -- shared look for the SoC figures (same palette as the slide decks and the
% core register placement figure). Every figure does % soc_style.tex -- shared look for the SoC figures (same palette as the slide decks and the
% core register placement figure). Every figure does % soc_style.tex -- shared look for the SoC figures (same palette as the slide decks and the
% core register placement figure). Every figure does \input{soc_style}.
\usepackage[T1]{fontenc}
\usepackage[scaled=0.95]{helvet}
\renewcommand{\familydefault}{\sfdefault}
\usepackage{booktabs}
\usepackage{amsmath}
\usepackage{amssymb}
\usetikzlibrary{arrows.meta,calc,positioning,fit,shapes.geometric,shapes.misc}

% ---------------- palette ----------------
\definecolor{pcb}{HTML}{0B2E33}
\definecolor{pcbmid}{HTML}{1F5C61}
\definecolor{tint}{HTML}{E6F0EF}
\definecolor{copper}{HTML}{B5651D}
\definecolor{coppert}{HTML}{F6E8DC}
\definecolor{ink}{HTML}{1F2A2E}
\definecolor{muted}{HTML}{5B6B70}
\definecolor{hardf}{HTML}{DCE6F2}
\definecolor{harde}{HTML}{2F4A7A}
\definecolor{good}{HTML}{2E7D32}
\definecolor{bad}{HTML}{B3261E}
% one colour per kind of block: border (d) / fill (f)
\definecolor{cpud}{HTML}{2B5C9E} \definecolor{cpuf}{HTML}{D6E4F5}
\definecolor{gpud}{HTML}{2E7D32} \definecolor{gpuf}{HTML}{D5EDD0}
\definecolor{npud}{HTML}{0B2E33} \definecolor{npuf}{HTML}{BFE0DB}
\definecolor{memd}{HTML}{5E3A8C} \definecolor{memf}{HTML}{E4DAF3}
\definecolor{busd}{HTML}{B5651D} \definecolor{busf}{HTML}{F6E8DC}
\definecolor{prd}{HTML}{546E7A}  \definecolor{prf}{HTML}{E3EAED}
\definecolor{secd}{HTML}{B3261E} \definecolor{secf}{HTML}{F7D9D6}
\definecolor{pwrd}{HTML}{8A6800} \definecolor{pwrf}{HTML}{FFF0B8}
\definecolor{medd}{HTML}{8E2463} \definecolor{medf}{HTML}{F5D6E6}
\definecolor{skd}{HTML}{2E7D32}  \definecolor{skf}{HTML}{BFE3B4}   % skew buffers, as in the register figure

\tikzset{
  blk/.style 2 args={draw=#1, fill=#2, rounded corners=4pt, thick, align=center,
                     font=\scriptsize, inner sep=3pt},
  cpu/.style={blk={cpud}{cpuf}},  gpu/.style={blk={gpud}{gpuf}},
  npu/.style={blk={npud}{npuf}},  mem/.style={blk={memd}{memf}},
  bus/.style={blk={busd}{busf}},  per/.style={blk={prd}{prf}},
  sec/.style={blk={secd}{secf}},  pwr/.style={blk={pwrd}{pwrf}},
  med/.style={blk={medd}{medf}},
  comb/.style ={draw=copper, fill=coppert, rounded corners=2pt, thick, align=center,
                font=\scriptsize, inner sep=2pt},                       % combinational logic
  hard/.style ={draw=harde, fill=hardf, thick, dashed, align=center,
                font=\scriptsize, inner sep=2pt},                       % hard macro / off-chip
  sram/.style ={draw=memd, fill=memf, thick, align=center, font=\scriptsize,
                inner sep=2pt},                                         % SRAM array, not flip-flops
  ff/.style   ={draw=#1, fill=#1!30, line width=0.35pt},                % one flip-flop square
  ffgone/.style={draw=muted, fill=white, densely dashed, line width=0.35pt},
  panel/.style={draw=muted, rounded corners=6pt, line width=0.9pt},
  ptitle/.style={anchor=north west, font=\normalsize\bfseries, text=ink, inner sep=0pt},
  note/.style ={font=\scriptsize, text=ink, align=left, inner sep=0pt},
  lab/.style  ={font=\tiny, text=muted, inner sep=1pt},
  axi/.style  ={{Stealth[length=2.6mm]}-{Stealth[length=2.6mm]}, line width=2.2pt, copper},
  axi1/.style ={-{Stealth[length=2.6mm]}, line width=2.2pt, copper},
  apb/.style  ={{Stealth[length=1.8mm]}-{Stealth[length=1.8mm]}, line width=1pt, copper},
  apb1/.style ={-{Stealth[length=1.8mm]}, line width=1pt, copper},
  dat/.style  ={-{Stealth[length=1.8mm]}, thick, pcbmid},
  ctl/.style  ={-{Stealth[length=1.8mm]}, densely dashed, thick, muted},
  irq/.style  ={-{Stealth[length=2mm]}, densely dotted, line width=1.3pt, bad},
  num/.style  ={circle, fill=copper, text=white, font=\scriptsize\bfseries, inner sep=1pt,
                minimum size=0.46cm},
  ours/.style ={draw=bad, line width=1.2pt, densely dashed, rounded corners=4pt},
}

% ---------------- helpers ----------------
\newcommand{\rn}[1]{\texttt{#1}}
% \ffrow{x}{y}{n}{cell}{style}: n flip-flops in a row, bottom-left corner at (x,y)
\newcommand{\ffrow}[5]{\foreach \q in {1,...,#3}{\draw[#5] ({#1+(\q-1)*#4},{#2}) rectangle ++(#4,#4);}}
% \ffcol{x}{y}{n}{cell}{style}: n flip-flops in a column, top-left corner at (x,y)
\newcommand{\ffcol}[5]{\foreach \q in {1,...,#3}{\draw[#5] ({#1},{#2-\q*#4}) rectangle ++(#4,#4);}}
% \ffgrid{x}{ytop}{n}{per row}{cell}{style}: n flip-flops, filled row by row from the top-left
\newcommand{\ffgrid}[6]{\pgfmathtruncatemacro{\gl}{#3-1}%
  \foreach \b in {0,...,\gl}{\pgfmathtruncatemacro{\gc}{mod(\b,#4)}\pgfmathtruncatemacro{\gr}{div(\b,#4)}%
    \draw[#6] ({#1+\gc*#5},{#2-(\gr+1)*#5}) rectangle ++(#5,#5);}}

% every figure: same page size, title, footer
\newcommand{\canvas}{\useasboundingbox (-0.3,-0.2) rectangle (38.6,30.4);}
\newcommand{\figtitle}[3]{%
  \node[num, minimum size=0.9cm, font=\Large\bfseries] at (0.45,29.55) {#1};
  \node[anchor=west, font=\LARGE\bfseries, text=pcb] at (1.15,29.55) {#2};
  \node[anchor=west, font=\small, text=ink] at (0,28.7) {#3};}
\newcommand{\figfoot}[1]{\node[anchor=west, font=\scriptsize, text=muted] at (0,0.25) {#1};}
% \panel{x0}{y0}{x1}{y1}{title}
\newcommand{\panel}[5]{\draw[panel] (#1,#2) rectangle (#3,#4);
  \node[ptitle] at ({#1+0.3},{#4-0.25}) {#5};}

% ---------------- timing diagrams: 1 cycle = \cw cm, low at y, high at y+0.45 ----------------
\newcommand{\cw}{1.0}
\newcommand{\wclk}[3]{% x0 y ncycles
  \foreach \k in {1,...,#3}{\draw[pcb, thick] ({#1+(\k-1)*\cw},#2) -- ++(0,0.45) -- ++(0.5*\cw,0)
     -- ++(0,-0.45) -- ++(0.5*\cw,0);}}
\newcommand{\wsig}[3]{% x0 y {level per cycle, e.g. 0,1,1,0}
  \gdef\wprev{x}%
  \foreach \v [count=\k from 0] in {#3}{%
    \draw[pcb, thick] ({#1+\k*\cw},{#2+0.45*\v}) -- ++(\cw,0);
    \if x\wprev\else\ifnum\wprev=\v\else\draw[pcb, thick] ({#1+\k*\cw},#2) -- ++(0,0.45);\fi\fi
    \global\let\wprev\v}}
\newcommand{\wbus}[5]{% x0 y c0 c1 text: one value held from cycle c0 to cycle c1
  \draw[pcbmid, thick, fill=white] ({#1+(#3)*\cw},{#2+0.225}) -- ({#1+(#3)*\cw+0.08},{#2+0.45})
    -- ({#1+(#4)*\cw-0.08},{#2+0.45}) -- ({#1+(#4)*\cw},{#2+0.225}) -- ({#1+(#4)*\cw-0.08},#2)
    -- ({#1+(#3)*\cw+0.08},#2) -- cycle;
  \node[font=\tiny\ttfamily] at ({#1+((#3)+(#4))*\cw/2},{#2+0.225}) {#5};}
\newcommand{\widle}[4]{% x0 y c0 c1: bus idle (no value)
  \draw[muted, thick] ({#1+(#3)*\cw},{#2+0.225}) -- ({#1+(#4)*\cw},{#2+0.225});}
\newcommand{\wname}[3]{\node[anchor=east, font=\scriptsize\ttfamily] at ({#1-0.15},{#2+0.225}) {#3};}
.
\usepackage[T1]{fontenc}
\usepackage[scaled=0.95]{helvet}
\renewcommand{\familydefault}{\sfdefault}
\usepackage{booktabs}
\usepackage{amsmath}
\usepackage{amssymb}
\usetikzlibrary{arrows.meta,calc,positioning,fit,shapes.geometric,shapes.misc}

% ---------------- palette ----------------
\definecolor{pcb}{HTML}{0B2E33}
\definecolor{pcbmid}{HTML}{1F5C61}
\definecolor{tint}{HTML}{E6F0EF}
\definecolor{copper}{HTML}{B5651D}
\definecolor{coppert}{HTML}{F6E8DC}
\definecolor{ink}{HTML}{1F2A2E}
\definecolor{muted}{HTML}{5B6B70}
\definecolor{hardf}{HTML}{DCE6F2}
\definecolor{harde}{HTML}{2F4A7A}
\definecolor{good}{HTML}{2E7D32}
\definecolor{bad}{HTML}{B3261E}
% one colour per kind of block: border (d) / fill (f)
\definecolor{cpud}{HTML}{2B5C9E} \definecolor{cpuf}{HTML}{D6E4F5}
\definecolor{gpud}{HTML}{2E7D32} \definecolor{gpuf}{HTML}{D5EDD0}
\definecolor{npud}{HTML}{0B2E33} \definecolor{npuf}{HTML}{BFE0DB}
\definecolor{memd}{HTML}{5E3A8C} \definecolor{memf}{HTML}{E4DAF3}
\definecolor{busd}{HTML}{B5651D} \definecolor{busf}{HTML}{F6E8DC}
\definecolor{prd}{HTML}{546E7A}  \definecolor{prf}{HTML}{E3EAED}
\definecolor{secd}{HTML}{B3261E} \definecolor{secf}{HTML}{F7D9D6}
\definecolor{pwrd}{HTML}{8A6800} \definecolor{pwrf}{HTML}{FFF0B8}
\definecolor{medd}{HTML}{8E2463} \definecolor{medf}{HTML}{F5D6E6}
\definecolor{skd}{HTML}{2E7D32}  \definecolor{skf}{HTML}{BFE3B4}   % skew buffers, as in the register figure

\tikzset{
  blk/.style 2 args={draw=#1, fill=#2, rounded corners=4pt, thick, align=center,
                     font=\scriptsize, inner sep=3pt},
  cpu/.style={blk={cpud}{cpuf}},  gpu/.style={blk={gpud}{gpuf}},
  npu/.style={blk={npud}{npuf}},  mem/.style={blk={memd}{memf}},
  bus/.style={blk={busd}{busf}},  per/.style={blk={prd}{prf}},
  sec/.style={blk={secd}{secf}},  pwr/.style={blk={pwrd}{pwrf}},
  med/.style={blk={medd}{medf}},
  comb/.style ={draw=copper, fill=coppert, rounded corners=2pt, thick, align=center,
                font=\scriptsize, inner sep=2pt},                       % combinational logic
  hard/.style ={draw=harde, fill=hardf, thick, dashed, align=center,
                font=\scriptsize, inner sep=2pt},                       % hard macro / off-chip
  sram/.style ={draw=memd, fill=memf, thick, align=center, font=\scriptsize,
                inner sep=2pt},                                         % SRAM array, not flip-flops
  ff/.style   ={draw=#1, fill=#1!30, line width=0.35pt},                % one flip-flop square
  ffgone/.style={draw=muted, fill=white, densely dashed, line width=0.35pt},
  panel/.style={draw=muted, rounded corners=6pt, line width=0.9pt},
  ptitle/.style={anchor=north west, font=\normalsize\bfseries, text=ink, inner sep=0pt},
  note/.style ={font=\scriptsize, text=ink, align=left, inner sep=0pt},
  lab/.style  ={font=\tiny, text=muted, inner sep=1pt},
  axi/.style  ={{Stealth[length=2.6mm]}-{Stealth[length=2.6mm]}, line width=2.2pt, copper},
  axi1/.style ={-{Stealth[length=2.6mm]}, line width=2.2pt, copper},
  apb/.style  ={{Stealth[length=1.8mm]}-{Stealth[length=1.8mm]}, line width=1pt, copper},
  apb1/.style ={-{Stealth[length=1.8mm]}, line width=1pt, copper},
  dat/.style  ={-{Stealth[length=1.8mm]}, thick, pcbmid},
  ctl/.style  ={-{Stealth[length=1.8mm]}, densely dashed, thick, muted},
  irq/.style  ={-{Stealth[length=2mm]}, densely dotted, line width=1.3pt, bad},
  num/.style  ={circle, fill=copper, text=white, font=\scriptsize\bfseries, inner sep=1pt,
                minimum size=0.46cm},
  ours/.style ={draw=bad, line width=1.2pt, densely dashed, rounded corners=4pt},
}

% ---------------- helpers ----------------
\newcommand{\rn}[1]{\texttt{#1}}
% \ffrow{x}{y}{n}{cell}{style}: n flip-flops in a row, bottom-left corner at (x,y)
\newcommand{\ffrow}[5]{\foreach \q in {1,...,#3}{\draw[#5] ({#1+(\q-1)*#4},{#2}) rectangle ++(#4,#4);}}
% \ffcol{x}{y}{n}{cell}{style}: n flip-flops in a column, top-left corner at (x,y)
\newcommand{\ffcol}[5]{\foreach \q in {1,...,#3}{\draw[#5] ({#1},{#2-\q*#4}) rectangle ++(#4,#4);}}
% \ffgrid{x}{ytop}{n}{per row}{cell}{style}: n flip-flops, filled row by row from the top-left
\newcommand{\ffgrid}[6]{\pgfmathtruncatemacro{\gl}{#3-1}%
  \foreach \b in {0,...,\gl}{\pgfmathtruncatemacro{\gc}{mod(\b,#4)}\pgfmathtruncatemacro{\gr}{div(\b,#4)}%
    \draw[#6] ({#1+\gc*#5},{#2-(\gr+1)*#5}) rectangle ++(#5,#5);}}

% every figure: same page size, title, footer
\newcommand{\canvas}{\useasboundingbox (-0.3,-0.2) rectangle (38.6,30.4);}
\newcommand{\figtitle}[3]{%
  \node[num, minimum size=0.9cm, font=\Large\bfseries] at (0.45,29.55) {#1};
  \node[anchor=west, font=\LARGE\bfseries, text=pcb] at (1.15,29.55) {#2};
  \node[anchor=west, font=\small, text=ink] at (0,28.7) {#3};}
\newcommand{\figfoot}[1]{\node[anchor=west, font=\scriptsize, text=muted] at (0,0.25) {#1};}
% \panel{x0}{y0}{x1}{y1}{title}
\newcommand{\panel}[5]{\draw[panel] (#1,#2) rectangle (#3,#4);
  \node[ptitle] at ({#1+0.3},{#4-0.25}) {#5};}

% ---------------- timing diagrams: 1 cycle = \cw cm, low at y, high at y+0.45 ----------------
\newcommand{\cw}{1.0}
\newcommand{\wclk}[3]{% x0 y ncycles
  \foreach \k in {1,...,#3}{\draw[pcb, thick] ({#1+(\k-1)*\cw},#2) -- ++(0,0.45) -- ++(0.5*\cw,0)
     -- ++(0,-0.45) -- ++(0.5*\cw,0);}}
\newcommand{\wsig}[3]{% x0 y {level per cycle, e.g. 0,1,1,0}
  \gdef\wprev{x}%
  \foreach \v [count=\k from 0] in {#3}{%
    \draw[pcb, thick] ({#1+\k*\cw},{#2+0.45*\v}) -- ++(\cw,0);
    \if x\wprev\else\ifnum\wprev=\v\else\draw[pcb, thick] ({#1+\k*\cw},#2) -- ++(0,0.45);\fi\fi
    \global\let\wprev\v}}
\newcommand{\wbus}[5]{% x0 y c0 c1 text: one value held from cycle c0 to cycle c1
  \draw[pcbmid, thick, fill=white] ({#1+(#3)*\cw},{#2+0.225}) -- ({#1+(#3)*\cw+0.08},{#2+0.45})
    -- ({#1+(#4)*\cw-0.08},{#2+0.45}) -- ({#1+(#4)*\cw},{#2+0.225}) -- ({#1+(#4)*\cw-0.08},#2)
    -- ({#1+(#3)*\cw+0.08},#2) -- cycle;
  \node[font=\tiny\ttfamily] at ({#1+((#3)+(#4))*\cw/2},{#2+0.225}) {#5};}
\newcommand{\widle}[4]{% x0 y c0 c1: bus idle (no value)
  \draw[muted, thick] ({#1+(#3)*\cw},{#2+0.225}) -- ({#1+(#4)*\cw},{#2+0.225});}
\newcommand{\wname}[3]{\node[anchor=east, font=\scriptsize\ttfamily] at ({#1-0.15},{#2+0.225}) {#3};}
.
\usepackage[T1]{fontenc}
\usepackage[scaled=0.95]{helvet}
\renewcommand{\familydefault}{\sfdefault}
\usepackage{booktabs}
\usepackage{amsmath}
\usepackage{amssymb}
\usetikzlibrary{arrows.meta,calc,positioning,fit,shapes.geometric,shapes.misc}

% ---------------- palette ----------------
\definecolor{pcb}{HTML}{0B2E33}
\definecolor{pcbmid}{HTML}{1F5C61}
\definecolor{tint}{HTML}{E6F0EF}
\definecolor{copper}{HTML}{B5651D}
\definecolor{coppert}{HTML}{F6E8DC}
\definecolor{ink}{HTML}{1F2A2E}
\definecolor{muted}{HTML}{5B6B70}
\definecolor{hardf}{HTML}{DCE6F2}
\definecolor{harde}{HTML}{2F4A7A}
\definecolor{good}{HTML}{2E7D32}
\definecolor{bad}{HTML}{B3261E}
% one colour per kind of block: border (d) / fill (f)
\definecolor{cpud}{HTML}{2B5C9E} \definecolor{cpuf}{HTML}{D6E4F5}
\definecolor{gpud}{HTML}{2E7D32} \definecolor{gpuf}{HTML}{D5EDD0}
\definecolor{npud}{HTML}{0B2E33} \definecolor{npuf}{HTML}{BFE0DB}
\definecolor{memd}{HTML}{5E3A8C} \definecolor{memf}{HTML}{E4DAF3}
\definecolor{busd}{HTML}{B5651D} \definecolor{busf}{HTML}{F6E8DC}
\definecolor{prd}{HTML}{546E7A}  \definecolor{prf}{HTML}{E3EAED}
\definecolor{secd}{HTML}{B3261E} \definecolor{secf}{HTML}{F7D9D6}
\definecolor{pwrd}{HTML}{8A6800} \definecolor{pwrf}{HTML}{FFF0B8}
\definecolor{medd}{HTML}{8E2463} \definecolor{medf}{HTML}{F5D6E6}
\definecolor{skd}{HTML}{2E7D32}  \definecolor{skf}{HTML}{BFE3B4}   % skew buffers, as in the register figure

\tikzset{
  blk/.style 2 args={draw=#1, fill=#2, rounded corners=4pt, thick, align=center,
                     font=\scriptsize, inner sep=3pt},
  cpu/.style={blk={cpud}{cpuf}},  gpu/.style={blk={gpud}{gpuf}},
  npu/.style={blk={npud}{npuf}},  mem/.style={blk={memd}{memf}},
  bus/.style={blk={busd}{busf}},  per/.style={blk={prd}{prf}},
  sec/.style={blk={secd}{secf}},  pwr/.style={blk={pwrd}{pwrf}},
  med/.style={blk={medd}{medf}},
  comb/.style ={draw=copper, fill=coppert, rounded corners=2pt, thick, align=center,
                font=\scriptsize, inner sep=2pt},                       % combinational logic
  hard/.style ={draw=harde, fill=hardf, thick, dashed, align=center,
                font=\scriptsize, inner sep=2pt},                       % hard macro / off-chip
  sram/.style ={draw=memd, fill=memf, thick, align=center, font=\scriptsize,
                inner sep=2pt},                                         % SRAM array, not flip-flops
  ff/.style   ={draw=#1, fill=#1!30, line width=0.35pt},                % one flip-flop square
  ffgone/.style={draw=muted, fill=white, densely dashed, line width=0.35pt},
  panel/.style={draw=muted, rounded corners=6pt, line width=0.9pt},
  ptitle/.style={anchor=north west, font=\normalsize\bfseries, text=ink, inner sep=0pt},
  note/.style ={font=\scriptsize, text=ink, align=left, inner sep=0pt},
  lab/.style  ={font=\tiny, text=muted, inner sep=1pt},
  axi/.style  ={{Stealth[length=2.6mm]}-{Stealth[length=2.6mm]}, line width=2.2pt, copper},
  axi1/.style ={-{Stealth[length=2.6mm]}, line width=2.2pt, copper},
  apb/.style  ={{Stealth[length=1.8mm]}-{Stealth[length=1.8mm]}, line width=1pt, copper},
  apb1/.style ={-{Stealth[length=1.8mm]}, line width=1pt, copper},
  dat/.style  ={-{Stealth[length=1.8mm]}, thick, pcbmid},
  ctl/.style  ={-{Stealth[length=1.8mm]}, densely dashed, thick, muted},
  irq/.style  ={-{Stealth[length=2mm]}, densely dotted, line width=1.3pt, bad},
  num/.style  ={circle, fill=copper, text=white, font=\scriptsize\bfseries, inner sep=1pt,
                minimum size=0.46cm},
  ours/.style ={draw=bad, line width=1.2pt, densely dashed, rounded corners=4pt},
}

% ---------------- helpers ----------------
\newcommand{\rn}[1]{\texttt{#1}}
% \ffrow{x}{y}{n}{cell}{style}: n flip-flops in a row, bottom-left corner at (x,y)
\newcommand{\ffrow}[5]{\foreach \q in {1,...,#3}{\draw[#5] ({#1+(\q-1)*#4},{#2}) rectangle ++(#4,#4);}}
% \ffcol{x}{y}{n}{cell}{style}: n flip-flops in a column, top-left corner at (x,y)
\newcommand{\ffcol}[5]{\foreach \q in {1,...,#3}{\draw[#5] ({#1},{#2-\q*#4}) rectangle ++(#4,#4);}}
% \ffgrid{x}{ytop}{n}{per row}{cell}{style}: n flip-flops, filled row by row from the top-left
\newcommand{\ffgrid}[6]{\pgfmathtruncatemacro{\gl}{#3-1}%
  \foreach \b in {0,...,\gl}{\pgfmathtruncatemacro{\gc}{mod(\b,#4)}\pgfmathtruncatemacro{\gr}{div(\b,#4)}%
    \draw[#6] ({#1+\gc*#5},{#2-(\gr+1)*#5}) rectangle ++(#5,#5);}}

% every figure: same page size, title, footer
\newcommand{\canvas}{\useasboundingbox (-0.3,-0.2) rectangle (38.6,30.4);}
\newcommand{\figtitle}[3]{%
  \node[num, minimum size=0.9cm, font=\Large\bfseries] at (0.45,29.55) {#1};
  \node[anchor=west, font=\LARGE\bfseries, text=pcb] at (1.15,29.55) {#2};
  \node[anchor=west, font=\small, text=ink] at (0,28.7) {#3};}
\newcommand{\figfoot}[1]{\node[anchor=west, font=\scriptsize, text=muted] at (0,0.25) {#1};}
% \panel{x0}{y0}{x1}{y1}{title}
\newcommand{\panel}[5]{\draw[panel] (#1,#2) rectangle (#3,#4);
  \node[ptitle] at ({#1+0.3},{#4-0.25}) {#5};}

% ---------------- timing diagrams: 1 cycle = \cw cm, low at y, high at y+0.45 ----------------
\newcommand{\cw}{1.0}
\newcommand{\wclk}[3]{% x0 y ncycles
  \foreach \k in {1,...,#3}{\draw[pcb, thick] ({#1+(\k-1)*\cw},#2) -- ++(0,0.45) -- ++(0.5*\cw,0)
     -- ++(0,-0.45) -- ++(0.5*\cw,0);}}
\newcommand{\wsig}[3]{% x0 y {level per cycle, e.g. 0,1,1,0}
  \gdef\wprev{x}%
  \foreach \v [count=\k from 0] in {#3}{%
    \draw[pcb, thick] ({#1+\k*\cw},{#2+0.45*\v}) -- ++(\cw,0);
    \if x\wprev\else\ifnum\wprev=\v\else\draw[pcb, thick] ({#1+\k*\cw},#2) -- ++(0,0.45);\fi\fi
    \global\let\wprev\v}}
\newcommand{\wbus}[5]{% x0 y c0 c1 text: one value held from cycle c0 to cycle c1
  \draw[pcbmid, thick, fill=white] ({#1+(#3)*\cw},{#2+0.225}) -- ({#1+(#3)*\cw+0.08},{#2+0.45})
    -- ({#1+(#4)*\cw-0.08},{#2+0.45}) -- ({#1+(#4)*\cw},{#2+0.225}) -- ({#1+(#4)*\cw-0.08},#2)
    -- ({#1+(#3)*\cw+0.08},#2) -- cycle;
  \node[font=\tiny\ttfamily] at ({#1+((#3)+(#4))*\cw/2},{#2+0.225}) {#5};}
\newcommand{\widle}[4]{% x0 y c0 c1: bus idle (no value)
  \draw[muted, thick] ({#1+(#3)*\cw},{#2+0.225}) -- ({#1+(#4)*\cw},{#2+0.225});}
\newcommand{\wname}[3]{\node[anchor=east, font=\scriptsize\ttfamily] at ({#1-0.15},{#2+0.225}) {#3};}

\newcommand{\tick}[2]{\node[circle, fill=good, text=white, font=\tiny\bfseries, inner sep=0.5pt] at (#1,#2) {\checkmark};}
\begin{document}
\begin{tikzpicture}[line join=round]
\canvas
\figtitle{6}{Interconnect: how blocks talk (AXI, AXI-Lite, APB, NoC)}
  {Every transfer is a handshake: the sender raises VALID, the receiver raises READY, and the data moves on the clock edge
   where both are 1. Your \rn{mm\_regs} bus already works this way.}

% =====================================================================
% (A) AXI4 channels
% =====================================================================
\panel{0}{13.6}{19.0}{28.1}{AXI4: five independent channels (32-bit address, 128-bit data)}
\node[npu, minimum width=2.6cm, minimum height=11.6cm, anchor=north west, text width=2.4cm] at (0.35,27.2)
  {\textbf{manager}\\(was ``master'')\\[4pt]e.g. NPU DMA};
\node[mem, minimum width=2.6cm, minimum height=11.6cm, anchor=north west, text width=2.4cm] at (16.05,27.2)
  {\textbf{subordinate}\\(was ``slave'')\\[4pt]e.g. DRAM controller};
\foreach \ch/\dir/\nm/\sigs/\wires [count=\k from 0] in {
    AW/1/write address/{AWID[3:0], AWADDR[31:0], AWLEN[7:0], AWSIZE[2:0], AWBURST[1:0], AWPROT[2:0]}/54,
    W/1/write data/{WDATA[127:0], WSTRB[15:0] (byte enables), WLAST}/147,
    B/0/write response/{BID[3:0], BRESP[1:0] (OKAY / error)}/8,
    AR/1/read address/{ARID[3:0], ARADDR[31:0], ARLEN[7:0], ARSIZE[2:0], ARBURST[1:0], ARPROT[2:0]}/54,
    R/0/read data/{RID[3:0], RDATA[127:0], RRESP[1:0], RLAST}/137}{
  \pgfmathsetmacro{\yy}{26.3-\k*2.3}
  \ifnum\dir=1 \draw[axi1] (3.0,\yy) -- (16.0,\yy); \else \draw[axi1] (16.0,\yy) -- (3.0,\yy); \fi
  \node[font=\scriptsize\bfseries, text=busd, anchor=south west, fill=white, inner sep=1pt] at (3.2,{\yy+0.12}) {\ch\ $\cdot$ \nm};
  \node[font=\scriptsize\bfseries, text=muted, anchor=south east, fill=white, inner sep=1pt] at (15.85,{\yy+0.12}) {\wires\ wires};
  \node[font=\tiny\ttfamily, anchor=north west, text width=12.6cm, align=left] at (3.2,{\yy-0.15}) {\sigs};
  \ifnum\dir=1
    \node[font=\tiny\ttfamily, text=bad, anchor=north west] at (3.2,{\yy-0.55}) {\ch VALID $\rightarrow$\quad $\leftarrow$ \ch READY};
  \else
    \node[font=\tiny\ttfamily, text=bad, anchor=north west] at (3.2,{\yy-0.55}) {\ch READY $\rightarrow$\quad $\leftarrow$ \ch VALID};
  \fi}
\node[note, anchor=north west, text width=18.2cm] at (0.35,15.3) {%
  Total: 54 + 147 + 8 + 54 + 137 = \textbf{400 wires} for one port (plus optional CACHE, LOCK, QOS).
  Reads and writes use separate channels, so both run at once. One address starts a burst of up to 256 beats
  (LEN + 1); IDs let many requests be in flight and come back out of order.};

% =====================================================================
% (B) handshake + burst waveforms
% =====================================================================
\panel{19.3}{13.6}{38.3}{28.1}{The VALID / READY rule, and a 4-beat read burst}
\renewcommand{\cw}{1.2}
\foreach \c in {0,...,11}{\node[lab] at ({22.6+(\c+0.5)*\cw},27.05) {\c};}
\wname{22.6}{26.35}{clk}     \wclk{22.6}{26.35}{12}
\wname{22.6}{25.55}{VALID}   \wsig{22.6}{25.55}{0,1,1,1,0,0,0,1,0,1,0,0}
\wname{22.6}{24.75}{READY}   \wsig{22.6}{24.75}{0,0,0,1,0,1,1,1,0,1,0,0}
\wname{22.6}{23.95}{DATA}
\widle{22.6}{23.95}{0}{1} \wbus{22.6}{23.95}{1}{4}{D0} \widle{22.6}{23.95}{4}{7}
\wbus{22.6}{23.95}{7}{8}{D1} \widle{22.6}{23.95}{8}{9} \wbus{22.6}{23.95}{9}{10}{D2} \widle{22.6}{23.95}{10}{12}
\foreach \c in {4,8,10}{\draw[good, thick, densely dashed] ({22.6+\c*\cw},23.8) -- ({22.6+\c*\cw},26.9); \tick{{22.6+\c*\cw}}{23.55}}
\node[lab, text=ink, anchor=north, text width=3.4cm, align=center] at ({22.6+2.0*\cw},23.4) {VALID first; waits for READY; D0 held};
\node[lab, text=ink, anchor=north, text width=3.0cm, align=center] at ({22.6+6.3*\cw},23.4) {READY first; data moves when VALID comes};
\node[lab, text=ink, anchor=north, text width=2.2cm, align=center] at ({22.6+10.2*\cw},23.4) {both at once};

\foreach \c in {0,...,11}{\node[lab] at ({22.6+(\c+0.5)*\cw},21.85) {\c};}
\wname{22.6}{21.15}{clk}      \wclk{22.6}{21.15}{12}
\wname{22.6}{20.4}{ARVALID}   \wsig{22.6}{20.4}{0,1,1,0,0,0,0,0,0,0,0,0}
\wname{22.6}{19.65}{ARREADY}  \wsig{22.6}{19.65}{0,0,1,0,0,0,0,0,0,0,0,0}
\wname{22.6}{18.9}{ARADDR}
\widle{22.6}{18.9}{0}{1} \wbus{22.6}{18.9}{1}{3}{0x8000\_0000, LEN=3} \widle{22.6}{18.9}{3}{12}
\wname{22.6}{18.15}{RVALID}   \wsig{22.6}{18.15}{0,0,0,0,0,1,1,1,1,1,0,0}
\wname{22.6}{17.4}{RREADY}    \wsig{22.6}{17.4}{1,1,1,1,1,1,1,0,1,1,1,1}
\wname{22.6}{16.65}{RDATA}
\widle{22.6}{16.65}{0}{5} \wbus{22.6}{16.65}{5}{6}{D0} \wbus{22.6}{16.65}{6}{7}{D1}
\wbus{22.6}{16.65}{7}{9}{D2 (held)} \wbus{22.6}{16.65}{9}{10}{D3} \widle{22.6}{16.65}{10}{12}
\wname{22.6}{15.9}{RLAST}     \wsig{22.6}{15.9}{0,0,0,0,0,0,0,0,0,1,0,0}
\foreach \c in {3,6,7,9,10}{\draw[good, thick, densely dashed] ({22.6+\c*\cw},15.85) -- ({22.6+\c*\cw},21.7);}
\node[note, anchor=north west, text width=18.2cm] at (19.6,15.55) {%
  Address accepted at the end of cycle 2. Data comes back after the memory latency. In cycle 7 the receiver is busy
  (RREADY = 0), so D2 stays on the bus one extra cycle: this is \textbf{back-pressure}. RLAST marks beat 4 of 4.
  Each green line = one transfer. Rule: VALID must never wait for READY, otherwise both sides could wait forever.};

% =====================================================================
% (C) protocol family
% =====================================================================
\panel{0}{0.7}{19.0}{13.3}{The bus family: which one where}
\node[anchor=north west, font=\scriptsize, inner sep=0pt] at (0.3,12.55) {%
  \renewcommand{\arraystretch}{1.2}%
  \begin{tabular}{@{}p{2.3cm} p{4.6cm} p{10.8cm}@{}}
    \toprule
    \textbf{Bus} & \textbf{Used for} & \textbf{What is special}\\
    \midrule
    AXI4        & CPU, GPU, NPU, DMA $\leftrightarrow$ memory & 5 channels, bursts up to 256 beats, IDs, many requests in flight\\
    AXI4-Lite   & control registers                           & 1 beat per request, 32 or 64 bit, no IDs: simple to build\\
    AXI4-Stream & data streams (video, DSP chains)            & no address at all: TVALID, TREADY, TDATA, TLAST\\
    APB         & slow peripherals (UART, I2C, timers)        & 2 phases (setup, access), one transfer at a time, very few gates (Figure 11)\\
    ACE / CHI   & CPU clusters, coherent caches               & ACE adds snoop channels to AXI; CHI is packet based for big NoCs\\
    Avalon-MM   & Intel FPGA (Platform Designer)              & read, write, waitrequest, readdatavalid (Figure 12)\\
    \rn{mm\_regs} bus & your design                           & valid / ready, write data sent with the address, read data 1 cycle later\\
                &                                             & $\Rightarrow$ like AXI-Lite with AW and W merged into one channel\\
    \bottomrule
  \end{tabular}};
\node[note, anchor=north west, text width=18.2cm] at (0.3,6.9) {%
  \textbf{Bridges} connect the families: an AXI $\rightarrow$ APB bridge lets the fast bus reach slow peripherals.
  Your UART bridge is exactly this kind of block: it turns UART bytes into \rn{mm\_regs} bus requests.\\[4pt]
  \textbf{Example system address map} (decided by the chip designers):\\[2pt]
  \rn{0x0000\_0000} boot ROM $\cdot$ \rn{0x0010\_0000} on-chip SRAM $\cdot$ \rn{0x1000\_0000} peripherals (APB)
  $\cdot$ \rn{0x2000\_0000} NPU registers $\cdot$ \rn{0x8000\_0000} DRAM.\\[2pt]
  Your own map is the same idea in small: \rn{0x0000} registers, \rn{0x1000} A, \rn{0x2000} B, \rn{0x3000} C.};

% =====================================================================
% (D) crossbar
% =====================================================================
\panel{19.3}{0.7}{28.6}{13.3}{Crossbar: 3 managers $\times$ 3 subordinates}
\foreach \nm [count=\k from 0] in {CPU, GPU, NPU}{
  \node[cpu, fill=white, minimum width=1.4cm, minimum height=0.8cm, font=\tiny\bfseries] (m\k) at (20.3,{10.9-\k*2.2}) {\nm};
  \node[comb, minimum width=1.2cm, minimum height=0.8cm, font=\tiny, text width=1.1cm] (d\k) at (22.1,{10.9-\k*2.2}) {address decode};
  \draw[dat] (m\k) -- (d\k);}
\foreach \nm [count=\k from 0] in {DRAM ctrl, SRAM, APB bridge}{
  \node[mem, fill=white, minimum width=1.6cm, minimum height=0.8cm, font=\tiny\bfseries, text width=1.4cm] (s\k) at (27.5,{10.9-\k*2.2}) {\nm};
  \node[comb, minimum width=1.25cm, minimum height=0.8cm, font=\tiny, text width=1.1cm] (a\k) at (25.6,{10.9-\k*2.2}) {arbiter + mux};
  \draw[dat] (a\k) -- (s\k);}
\foreach \i in {0,1,2}{\foreach \j in {0,1,2}{\draw[busd, thick] (d\i.east) -- (a\j.west);}}
\node[note, font=\tiny, anchor=north west, text width=8.9cm] at (19.6,5.95) {%
  The \textbf{decoder} compares the address with the map and picks the path. If two managers want the same subordinate
  in one cycle, its \textbf{arbiter} picks one (round-robin or by priority / QoS); the other waits (READY = 0).\\[3pt]
  All combinational logic plus a few registers for the arbiter state. Fine for a few blocks; for 20+ blocks the wires
  explode $\Rightarrow$ network-on-chip.};

% =====================================================================
% (E) NoC router
% =====================================================================
\panel{28.9}{0.7}{38.3}{13.3}{Network-on-chip: routers and packets}
\foreach \r in {0,1,2}{\foreach \c in {0,1,2}{
  \node[draw=busd, fill=busf, thick, minimum size=0.55cm, font=\tiny\bfseries] (n\r\c) at ({29.9+\c*1.3},{11.6-\r*1.3}) {R};}}
\foreach \r in {0,1,2}{\draw[busd, thick] (n\r0) -- (n\r1) -- (n\r2);}
\foreach \c in {0,1,2}{\draw[busd, thick] (n0\c) -- (n1\c) -- (n2\c);}
\node[note, font=\tiny, anchor=north west, text width=4.0cm] at (34.0,12.2) {%
  3$\times$3 mesh: each block plugs into one router. A packet hops router to router (1--2 clocks per hop).};
% router internals
\node[anchor=north west, font=\scriptsize\bfseries] at (29.2,8.3) {inside one router (5 ports: N, S, E, W, local)};
\foreach \p [count=\k from 0] in {N, S, E, W, local}{
  \node[font=\tiny\bfseries, anchor=east] at (29.95,{7.55-\k*0.62}) {\p};
  \ffrow{30.0}{7.4-\k*0.62}{4}{0.3}{ff=busd}}
\node[lab, anchor=north] at (30.6,4.4) {input FIFOs (flit buffers = FF)};
\node[comb, minimum width=1.5cm, minimum height=1.1cm, font=\tiny, text width=1.4cm] at (32.55,6.95) {route compute (XY)};
\node[comb, minimum width=1.5cm, minimum height=1.1cm, font=\tiny, text width=1.4cm] at (32.55,5.55) {switch arbiter};
\node[comb, minimum width=1.6cm, minimum height=2.6cm, font=\tiny, text width=1.4cm] at (34.55,6.25) {5$\times$5 crossbar};
\foreach \k in {0,...,4}{\ffrow{35.55}{7.4-\k*0.62}{1}{0.3}{ff=busd}}
\node[lab, anchor=north west, text width=2.2cm] at (35.95,7.65) {output register per port};
\node[note, font=\tiny, anchor=north west, text width=8.9cm] at (29.2,3.85) {%
  A packet is cut into \textbf{flits}: the head flit carries the destination, body flits carry data, the tail flit closes the path.
  XY routing: go along X until the column matches, then along Y. Same handshake idea per link (valid / credit).};

\figfoot{Signal names follow AMBA AXI4 / APB. Widths are one example configuration; the address map is an example.}
\end{tikzpicture}
\end{document}
